% theme: general_angles_and_radian_measure
\MajorQuestionLesson{一般角と弧度法}
\SetRowSpace{5mm}

\TypeQuestionDouble{次の角と同じ動径を表す角のうち、$0^\circ \leqq \theta < 360^\circ$ を満たす角 $\theta$ を求めなさい。}{kihon}
\QQRow{$390^\circ$}{}{$-45^\circ$}{}
\QQRow{$765^\circ$}{}{$-510^\circ$}{}
\QQRow{$1080^\circ$}{}{$-870^\circ$}{}
\QQRow{$225^\circ$}{}{$-270^\circ$}{}

\TypeQuestionDouble{次の角の動径は第何象限にあるか答えなさい。座標軸上にある場合は、どの座標軸の正または負の部分にあるか答えなさい。}{shougen}
\QQRow{$30^\circ$}{}{$135^\circ$}{}
\QQRow{$240^\circ$}{}{$330^\circ$}{}
\QQRow{$90^\circ$}{}{$-180^\circ$}{}
\QQRow{$750^\circ$}{}{$-585^\circ$}{}

\TypeQuestionDouble{次の角と同じ動径を表す一般角を、$0^\circ \leqq \alpha < 360^\circ$ として、$\alpha+360^\circ n$（$n$ は整数）の形で表しなさい。}{ippankaku}
\QQRow{$75^\circ$}{}{$-40^\circ$}{}
\QQRow{$150^\circ$}{}{$270^\circ$}{}
\QQRow{$420^\circ$}{}{$-225^\circ$}{}

\TypeQuestionDouble{次の2つの角が同じ動径を表す場合は○、表さない場合は×を答えなさい。}{doukei}
\QQRow{$30^\circ,\ 390^\circ$}{}{$-45^\circ,\ 315^\circ$}{}
\QQRow{$120^\circ,\ -240^\circ$}{}{$150^\circ,\ 510^\circ$}{}
\QQRow{$60^\circ,\ 300^\circ$}{}{$-90^\circ,\ 90^\circ$}{}
\QQRow{$720^\circ,\ 0^\circ$}{}{$-135^\circ,\ 225^\circ$}{}

\LessonPageBreak

\TypeQuestionSingle{次の角を弧度法で表しなさい。}{degree_to_radian}
\QQRow{$30^\circ$}{}{$45^\circ$}{}
\QQRow{$120^\circ$}{}{$225^\circ$}{}
\QQRow{$-150^\circ$}{}{$300^\circ$}{}
\QQRow{$540^\circ$}{}{$72^\circ$}{}

\TypeQuestionSingle{次の角を度数法で表しなさい。}{radian_to_degree}
\QQRow{$\dfrac{\pi}{6}$}{}{$\dfrac{5\pi}{4}$}{}
\QQRow{$-\dfrac{2\pi}{3}$}{}{$\dfrac{7\pi}{6}$}{}
\QQRow{$\dfrac{11\pi}{3}$}{}{$\dfrac{3\pi}{10}$}{}
\QQRow{$-\dfrac{7\pi}{4}$}{}{$\dfrac{13\pi}{6}$}{}

\TypeQuestionDouble{次の角と同じ動径を表す一般角を、$0 \leqq \alpha < 2\pi$ として、$\alpha+2n\pi$（$n$ は整数）の形で表しなさい。}{ippankaku_rad}
\QQRow{$\dfrac{\pi}{4}$}{}{$-\dfrac{\pi}{3}$}{}
\QQRow{$\dfrac{13\pi}{6}$}{}{$-\dfrac{9\pi}{4}$}{}

\TypeQuestionDouble{半径 $r$ の円において、長さ $l$ の弧に対する中心角を弧度法で求めなさい。}{radian_meaning}
\QQRow{$r=6,\ l=9$}{}{$r=10,\ l=4\pi$}{}
\QQRow{$r=3,\ l=2$}{}{$r=8,\ l=12$}{}

\TypeQuestionDouble{半径と中心角が次のように与えられた扇形について、弧の長さと面積を求めなさい。中心角は弧度法で表されているものとします。}{ogi}
\QQ{半径 $6$、中心角 $\dfrac{\pi}{3}$}{}
\QQ{半径 $4$、中心角 $\dfrac{3\pi}{4}$}{}
\QQ{半径 $9$、中心角 $\dfrac{2\pi}{3}$}{}
\QQ{半径 $5$、中心角 $2$}{}
